\documentclass[11pt]{amsart}

\usepackage[a4paper,left=2.7cm,right=2.7cm,top=2.7cm,bottom=2.7cm]{geometry}
\usepackage{amsmath,amssymb,amsthm}
\usepackage[shortlabels]{enumitem}
\usepackage{booktabs}
\usepackage[dvipsnames]{xcolor}
\usepackage[colorlinks=true,citecolor=MidnightBlue,linkcolor=BrickRed,urlcolor=MidnightBlue]{hyperref}

\newtheorem{theorem}{Theorem}[section]
\newtheorem{proposition}[theorem]{Proposition}
\newtheorem{lemma}[theorem]{Lemma}
\newtheorem{corollary}[theorem]{Corollary}
\theoremstyle{definition}
\newtheorem{definition}[theorem]{Definition}
\newtheorem{remark}[theorem]{Remark}

\newcommand{\Pcal}{\mathcal P}
\newcommand{\Tcal}{\mathcal T}
\newcommand{\DHL}{\operatorname{DHL}}
\title[Bounded gaps between primes]{Bounded gaps between primes:\
towards the bound $H_1\leq236$}
\author{\NoCaseChange{Roy van Rijn}}
\date{Working draft, 3 September 2026}
\thanks{This is a conditional research draft.  It does not claim an unconditional
improvement until the analytic and computational obligations stated in the paper
have been discharged.}
\hypersetup{
  pdftitle={Bounded gaps between primes: towards the bound H1 <= 236},
  pdfauthor={Roy van Rijn},
  pdfsubject={A conditional route to the prime-gap bound 236},
  pdfkeywords={bounded prime gaps, Maynard--Tao sieve, equidistribution,
    exact certificate, Lean formalization}
}

\begin{document}

\begin{abstract}
We describe a route from the bound $H_1\leq240$ of Stadlmann to the next
admissible-tuple threshold $H_1\leq236$.  The argument uses the Maynard--Tao
sieve on a two-band support, an extension of the Type~IIc distribution range,
and a fixed degree-$27$ rational variational candidate in $48$ variables.  The
finite tuple calculation, the prime-gap endgame, the candidate metadata, and
the scalar certificate arithmetic have been checked in Lean against a pinned
revision of PrimeGapsLib.  At the time of this draft, the outward-rounded
boundary integral calculation and the final analytic review remain open.  We
state the remaining hypotheses separately so that the completed components can
be audited without treating numerical evidence as a theorem.
\end{abstract}

\maketitle

\section{Introduction}\label{sec:introduction}

Let $p_n$ denote the $n$-th prime and put
\[
  H_1=\liminf_{n\to\infty}(p_{n+1}-p_n).
\]
Zhang proved that $H_1$ is finite~\cite{Zhang2014}.  The Maynard--Tao sieve
subsequently led to much smaller bounds~\cite{Maynard2015}, and the Polymath8b
project obtained $H_1\leq246$~\cite{PolymathVSS}.  Stadlmann combined the
Bombieri--Vinogradov range with distribution estimates for moduli having a
large smooth factor and proved $H_1\leq240$~\cite{Stadlmann2026}.

The next finite threshold is supplied by an admissible $48$-tuple of diameter
$236$.  Thus it is enough to establish $\DHL[48,2]$: every admissible
$48$-tuple has infinitely many translates containing at least two primes.  The
purpose of this paper is to isolate a concrete route to that statement and to
make the computational and analytic proof boundaries explicit.

\begin{theorem}[Conditional assembly]\label{thm:conditional}
Assume the distribution estimate in Proposition~\ref{prop:distribution} and
the certified variational inequality in Proposition~\ref{prop:variational}.
Then
\[
  H_1\leq236.
\]
\end{theorem}

\paragraph{Current status.}
The hypotheses of Theorem~\ref{thm:conditional} are not both certified at
the date of this draft.  In particular, the theorem is not an unconditional
claim that $H_1\leq236$.

\subsection{Proof strategy}

We follow the organization of Stadlmann's proof while replacing the two pieces
that determine the numerical threshold.  First, we use the same shaped-support
version of the Maynard--Tao sieve and reduce $\DHL[48,2]$ to a distribution
statement and a strict inequality between two quadratic integrals.  Second, we
extend the terminal Type~IIc inequality just far enough to permit the rational
endpoint $A_2=2029/8000$.  Third, for one fixed rational polynomial $F$, we
evaluate the unrestricted part of the integrals exactly and enclose only the
support-boundary correction.  Finally, an explicit admissible tuple converts
$\DHL[48,2]$ to the asserted prime-gap bound.

The proof is arranged so that expensive computation is never implicit in a
test.  A source candidate, its exact unrestricted contractions, and the
outward-rounded boundary result are hash-identified artifacts.  Cheap replay
checks the hashes and exact inequalities.  The finite endgame is proved in
Lean using PrimeGapsLib~\cite{PrimeGapsLib}.

\subsection{Structure of the paper}

Section~\ref{sec:sieve} records the shaped-support sieve reduction.
Section~\ref{sec:distribution} states the distribution improvement and the
new Type~IIc terminal condition.  Section~\ref{sec:parameters} fixes the exact
parameters.  Section~\ref{sec:certificate} describes the rational variational
certificate.  Section~\ref{sec:endgame} gives the formalized finite endgame,
and Section~\ref{sec:assembly} lists the precise remaining certification steps.

\subsection{Notation}

Lower-case Roman letters generally denote integers, and $p$ denotes a prime.
We write $1_{\Pcal}$ for the prime indicator, $P(y)=\prod_{p<y}p$, and use
Vinogradov notation with subscripts to indicate permitted dependence.  All
rational endpoint comparisons in the final parameter and certificate checks
are meant literally, not as floating-point comparisons.

\section{The shaped-support sieve}\label{sec:sieve}

We recall only the portion of the Maynard--Stadlmann framework needed for the
present parameters.  Full sieve asymptotics may be taken from
\cite{Stadlmann2026}; their formalization is a separate part of the program
described below.

\begin{definition}[Two-band support]\label{def:support}
Let $k\geq1$, $\delta,\varepsilon>0$, and
\[
  -\varepsilon=A_0<A_1<A_2<\tfrac12-\varepsilon.
\]
For each band $j\in\{1,2\}$ let $(B_{j,m})_{m\geq1}$ be nondecreasing.  We
write $\Tcal_k(\delta,\mathbf A,\mathbf B,\varepsilon)$ for the union over
$j=1,2$ of the points $\mathbf t\in[0,1]^k$ satisfying
\[
 A_{j-1}+\varepsilon\leq \sum_{i=1}^k t_i<A_j+\varepsilon,
 \qquad
 \sum_{i:t_i>\delta}t_i\leq B_{j,m},
 \quad m=\#\{i:t_i>\delta\}.
\]
\end{definition}

Let $F$ be a symmetric square-integrable function supported on this region.
Write
\[
 I(F)=\int_{\Tcal_k}F(\mathbf t)^2\,d\mathbf t
\]
and let $J(F)$ denote the one-coordinate marginal quadratic integral appearing
in the Maynard--Tao sieve.  We use the normalization in which the required
strict inequality is
\begin{equation}\label{eq:variational-target}
  kJ(F)>I(F).
\end{equation}

\begin{proposition}[Shaped-support sieve criterion]\label{prop:sieve}
Suppose the prime indicator has the required equidistribution for every modulus
arising from a pair of points in
$\Tcal_k(\delta,\mathbf A,\mathbf B,\varepsilon)$.  If a symmetric $F$
supported on this region satisfies \eqref{eq:variational-target}, then
$\DHL[k,2]$ holds.
\end{proposition}

\begin{remark}
For the parameters below we use the prime indicator itself.  Consequently no
minorant-loss $K$ term occurs in \eqref{eq:variational-target}.  This is the
same simplification used in the final $240$ specialization of
\cite{Stadlmann2026}.
\end{remark}

\paragraph{Formalization boundary.}
Before unconditional promotion, Proposition~\ref{prop:sieve} must be
transcribed with all uniformity and smoothing hypotheses and connected to the
formal $k=48$ statement in Lean.

\section{The distribution estimate}\label{sec:distribution}

The only parameter face that prevents the published distribution criterion
from reaching the endpoint used here is the terminal Type~IIc estimate.  Put
\[
  \omega=A_2-\frac14,
  \qquad \gamma=\frac25.
\]
The incomplete-rectangle analysis leads to the following replacement for the
old terminal condition.

\begin{proposition}[Incomplete-rectangle Type~IIc estimate]
\label{prop:distribution}
For the originating Type~IIc family in the Harman decomposition, retain the
actual incomplete ranges, affine congruence, conductor reductions, and source
divisors.  If
\begin{equation}\label{eq:typeiic}
  6-22\gamma+72\delta+216\omega<0,
\end{equation}
then the Type~IIc contribution has the power saving required by the
distribution criterion for the support in Definition~\ref{def:support}.
Together with the unchanged Type~I, Type~II, and Type~III estimates, this gives
the equidistribution hypothesis of Proposition~\ref{prop:sieve}.
\end{proposition}

We indicate the source of the saving.  Write
\[
 G=(\lambda,\widetilde\lambda)=w_2g,
 \qquad \lambda=Ga,\quad\widetilde\lambda=Gb,
 \qquad (a,b)=1.
\]
For tuples inherited from the pre-completion sum, $G$ divides the shift
parameter and the lifted variables lie on an affine line
\[
  aY-bX=d\{cj+(a-b)B\}.
\]
For $g$ above the critical threshold, one-dimensional completion on this line
supplies the missing factor.  For small $g$, the two-dimensional Fourier slice
is a rank-three Kloosterman pullback.  Translated correlations are controlled
using geometric irreducibility and trace-function quasi-orthogonality
\cite{Katz1988,FKM2017}; a Plancherel average replaces the false pointwise
bound at zero frequency.

The conductor bookkeeping is part of the assertion.  In particular, the
A-process uses the source divisor $q_A$, not an unproved smooth factor of the
whole modulus.  The factors coming from $(q_0,t)$, $(q_0,s)$, and the
phase-dead part $(q_3,t)$ remain visible through the completion argument.

\paragraph{Verification boundary.}
A complete typeset proof of Proposition~\ref{prop:distribution}, including
the inactive parameter faces and all endpoint uniformity, must be independently
reviewed.  The repository contains checked derivations and hostile finite and
exponent audits, but those are not a substitute for this proof.

\section{Permissible parameters}\label{sec:parameters}

We fix
\begin{equation}\label{eq:parameters}
 k=48,\quad
 \delta=\frac7{250},\quad
 \varepsilon=\frac{17}{2000},\quad
 \mathbf A=\left(-\frac{17}{2000},\frac{23}{100},
                         \frac{2029}{8000}\right).
\end{equation}
Only the first seven large-coordinate counts can be active.  The two rows of
$\mathbf B$ are
\begin{align*}
 \mathbf B_1&=\left(\frac9{50},\frac9{50},
  \frac15,\frac15,\frac15,\frac15,\frac15\right),\\
 \mathbf B_2&=\left(\frac3{20},\frac3{20},
  \frac{17}{100},\frac{17}{100},\frac{17}{100},
  \frac{17}{100},\frac{17}{100}\right).
\end{align*}
The outer and shared-coordinate radii are respectively
\[
 A_2+\varepsilon=\frac{2097}{8000},
 \qquad A_2-\varepsilon=\frac{1961}{8000}.
\]

For these parameters, $\omega=29/8000$, and the left side of
\eqref{eq:typeiic} is exactly
\[
  6-22\left(\frac25\right)+72\left(\frac7{250}\right)
      +216\left(\frac{29}{8000}\right)=-\frac1{1000}.
\]
Thus the selected rational endpoint lies strictly inside the proposed
Type~IIc range.  The remaining distribution faces are checked by an
exact-rational finite oracle implementing the support-cell criteria from
\cite{Stadlmann2026}.

\paragraph{Matching requirement.}
The finite oracle and Proposition~\ref{prop:distribution} must be matched line
by line: the same strict epsilon convention, support rows, modulus
factorization, and minorant choice must occur in both.

\section{The variational certificate}\label{sec:certificate}

The candidate is a symmetric rational polynomial of degree at most $27$ in
$48$ variables, represented by $2526$ nonzero orbit coefficients.  Its
canonical source file has SHA-256 digest
\begin{center}
{\small\ttfamily d1cecf4db18f29da80cb2f5784ecfb9ee98ccc1221055b4a45ed3844427339ac}.
\end{center}
The coefficient table is supplied as ancillary material rather than printed in
the paper.

After a common positive normalization, exact rational contraction on the
outer simplex gives
\begin{align}
 I_{\rm legal}&\leq I_{\rm simplex}
      =0.9999955585079013\ldots,\label{eq:Iupper}\\
 48J_{\rm unrestricted}&=1.0005162874604419\ldots.
      \label{eq:Junrestricted}
\end{align}
The first inequality uses only containment of the legal support in the outer
simplex and nonnegativity of $F^2$.

Let $\Delta_B$ be the normalized legal-minus-unrestricted boundary correction
to $48J$.  Equations \eqref{eq:Iupper} and
\eqref{eq:Junrestricted} reduce the remaining numerical task to
\begin{equation}\label{eq:boundary-target}
  \Delta_B>-0.0005207289525405\ldots.
\end{equation}
No floating-point value is used to certify this inequality.  The implementation
partitions the boundary contribution into finitely many geometry cells,
evaluates every cell with outward-rounded Arb balls, and sums the resulting
rational endpoints.  The final lower endpoint is then compared with the exact
rational difference $I_{\rm simplex}-48J_{\rm unrestricted}$.

\begin{proposition}[Variational certificate]\label{prop:variational}
For the exact parameters \eqref{eq:parameters} and the rational candidate
identified above, the legal-support integrals satisfy
\[
  48J(F)>I(F).
\]
\end{proposition}

\paragraph{Current status.}
Proposition~\ref{prop:variational} remains pending until the full boundary run
is complete, its checkpoint and manifest hashes are frozen, and the
outward-rounding semantics are connected to the integral decomposition.  The
existing numerical corrections are discovery evidence only.

\section{The admissible tuple and formal endgame}\label{sec:endgame}

Consider the set
\begin{multline*}
\mathcal H=\{0,6,8,14,18,24,26,48,50,54,56,60,66,68,74,78,\\
80,84,90,96,98,104,110,116,120,126,134,138,144,150,158,164,\\
168,176,180,186,188,194,200,204,206,210,216,224,228,230,234,236\}.
\end{multline*}
It contains $48$ elements and has diameter $236$.  It is admissible: for every
prime $p\leq47$ an explicit residue class is omitted, while for $p>48$ the
cardinality bound is automatic.  Such tuples are tabulated in the narrow
admissible tuple data of Sutherland~\cite{Sutherland}.

The following implication is already machine checked.

\begin{proposition}[Formalized endgame]\label{prop:endgame}
If $\DHL[48,2]$ holds, then infinitely many consecutive prime gaps have size at
most $236$.
\end{proposition}

The Lean proof verifies the cardinality, diameter, and admissibility of
$\mathcal H$ by decision procedures and then applies the generic interval-to-
consecutive-gap argument from PrimeGapsLib.  The project pins the PrimeGapsLib
commit, Lean toolchain, Mathlib revision, and complete transitive dependency
manifest.  No expensive variational computation is run by the Lean build.

\section{Assembly and certification status}\label{sec:assembly}

Assuming Propositions~\ref{prop:distribution} and
\ref{prop:variational}, Proposition~\ref{prop:sieve} gives
$\DHL[48,2]$.  Proposition~\ref{prop:endgame} then proves
Theorem~\ref{thm:conditional}.

For clarity, the status of each logical layer is as follows.
\begin{center}
\begin{tabular}{@{}p{0.31\linewidth}p{0.57\linewidth}@{}}
\toprule
Layer & Status at this draft \\
\midrule
Admissible $48$-tuple and diameter $236$ & Checked independently and in Lean.\\
Exact unrestricted $I$ and $48J$ & Exact rational computations frozen and
replayed; scalar comparison formalized in Lean.\\
Boundary correction $\Delta_B$ & Full outward-rounded calculation in progress;
final certificate not yet claimed.\\
Type~IIc extension & Research proof and adversarial audits recorded; complete
typeset proof and independent human review pending.\\
Shaped-support sieve bridge & Mathematical criterion available from the
Stadlmann framework; exact specialization and Lean connection to $\DHL[48,2]$
pending.\\
\bottomrule
\end{tabular}
\end{center}

\section{Reproducibility}

The repository separates expensive computation from replay.  Candidate files,
checkpoints, manifests, and exact contractions are identified by SHA-256.  The
default tests verify schemas, hashes, exact rational identities, and theorem
interfaces without repeating the boundary integration.  The Lean target can be
built with
\begin{verbatim}
cd formalization
lake build PrimeGaps236
\end{verbatim}
After completion, the D27 finalizer emits an exact rational lower endpoint.  A
guarded exporter accepts only a complete, certification-eligible, hash-matched,
strictly crossing result and produces a small Lean theorem whose rational
comparison is proved by kernel-checked normalization.

\section*{Acknowledgements and declaration of generative AI use}

OpenAI Codex was used during literature search, computational experimentation,
consistency checking, Lean formalization, and preparation of this manuscript.
The author reviewed the mathematical claims and assumes full responsibility for
the contents.  Repository and archive identifiers will be added before
circulation.

\bibliographystyle{amsplain}
\bibliography{references}

\end{document}
